% !TeX root = ../papers/nc-p3.tex
\section{Pauli射影閉包の境界での正則性の完全判定}
\label{sec:pauli-projective-boundary}
\begin{theorem}\label{thm:pauli-projective-boundary}
標数零の代数閉体上で、規格化テンソルの五つの二次式が定める射影曲面を考える。
$c=a+b$ とする。開集合 $abcLM\ne0$ の境界において、反対階数と
その四次テンソルの巡回隣接像が与える窓の標準14641型正則性は次の通りである。
\[
\begin{array}{c|c|c|l}
 & LM &\text{反対階数}&\text{判定}\\\hline
 abc=0,\ (a,b)\ne(0,0)&\ne0&14&\text{全点正則}\\
 a=b=0 &\ne0&12&\text{一つの量子同型類}\\
 abc\ne0&0&9&\text{長さ4の次元が36以上}\\
 abc=0,\ (a,b)\ne(0,0)&0&4&\text{隣接階数4}
\end{array}
\]
正則な両層では四中心四次式の多項式環上で各始点の経路の直和が自由階数256になる。
\end{theorem}
\begin{proof}
反対階数は非零なPauli係数の個数に等しい。$LM=-AB=-CD$、
三つの平方和はそれぞれ $bc,ac,ab$ である。$a=b=LM=0$ は射影点を与えない。
従って支持を数えれば表の四層で尽きる。192変換は $a,b,c$ の各々を
単元倍を除いて新しい $a$ に取れる。三つのうちちょうど一つが零のとき、
零のものを新しい $a$ とすることも、非零のものを新しい $a$ とすることもできる。

階数14では $a=1$ と取ると $b=0,-1$、従って $K=1,-1$ である。
復元式 $s=B-A+D-C$, $t=B+A+D+C$, $n=2L$ とその逆元の式は
$U=1$ でも五つの二次式だけから従う。$H=4LM\ne0$、$n^4=-H^2$ である。
$K+H\ne0,\pm1$ かつ $K-H\ne\pm1$ なら
定理\ref{thm:rank-fifteen-covers}の通常な階数14窓である。
Pauli係数の恒等式はこの通常域でも成り立つ。実際Gram表示の変換と
四次根による係数式はいずれもここで分母が単元である。
正則性と五次元の四次中心は定理\ref{thm:rank-fourteen-all-ordinary}と
定理\ref{thm:rank-fourteen-center}による。

残るのは $H^2=1,4$ に限られる。
\path{computations/verify_pauli_projective_boundary.py} は101上で $K,H$ の両符号、
$r,w,s,t,n$ の全根から形式的なテンソルを直接作る。
根の選択は順に256通りと128通りで、規格化した192軌道では二代表と一代表になる。
全代表で四始点の双対次元 $(1,4,6,4,1,0)$、経路長4,5の次元 $35,56$、
中心の次元 $(3,2,0,0)$、四中心商の長さ13までの標準列を確認した。
標数零の各根はこれらのいずれかに還元される。
隣接階数6と一文字階数4（反対階数14より従う）から双対の完全対合を得る。
関係と中心の定階数による持上げ、一般点からの経路次元の下界と四中心商による上界から
全次数の主張が従う。階数12は定理\ref{thm:rank-twelve-quantum-limit}である。

階数9では $a=1$ とし、$P=(u+u^{-1})/2$, $Q=(u-u^{-1})/2$ と置く。
$b=Q^2,\ c=P^2$ で、$(L,M),(A,B),(C,D)$ は各対の一方だけが非零である。
その三つの非零成分は順に $\pm PQ,P,Q$ と取れる。実際後の二成分を
$P,Q$ として $P^2-Q^2=1$ なので $u=P+Q$ と復元できる。
非零成分の位置と初めの符号による十六表示の各々で、有理関数体上の長さ4の
全ての関係挿入を文字別に行消去した。始点0で得る商次元は
$(10,10,8,8)$ で、合計36である。特殊点では元のテンソルの生成行を用いる
全挿入行列の階数は増えない。従って全点で標準次元35を超える。
階数4では $a=0,b=1$ とすれば $L,M$ の一方のみが $\pm1$ である。
この四テンソルの隣接階数は四始点とも4である。

これらの境界層はいずれも高々一次元である。五つの斉次方程式の
$\mathbb P^7$ における各既約成分は次元2以上なので、境界に含まれる別成分はない。
従ってこれは稠密な形式的開集合の射影閉包の境界全体の判定である。
\end{proof}

\begin{corollary}\label{cor:pauli-rank-fourteen-valuation-wall}
剰余標数が零または101で $\nu(y)=2\nu(x)<0$ なら、標数零の通常な階数15窓は
標準Koszul左右AS正則で、各始点で四次中心式四本の上に自由階数256となる。
\end{corollary}
\begin{proof}
前節の $D=x^3/y$ は極で、$U$ が1に還元される根を取れる。
$V=(U+3)^2 y/x^2$ より $\bar V=16\,\overline{y/x^2}\ne0$ である。
順にモニック方程式の根を取れば、表示は階数14の形式的なテンソルに還元される。
剰余標数零では定理の証明から、各剰余窓の四次中心は五次元で
四中心商の標準列が得られる。101の $\bar V=1,4$ でも上の直接計算が適用できる。

101の通常な階数14還元では定理\ref{thm:rank-fourteen-geometric-reduction}の
多項式核計算が四中心商を与える。ここでは階数15方向に中心が持ち上がるよう、
剰余の全四次中心の次元も検査する。
\path{computations/verify_rank_fourteen_wall_centers.py} は既存の階数14の
関数体上で、長さ5までの各経路行列と四つの中心方程式の主元の積を取る。
これは選択した行と列の実際の小行列式である。ノルムの分子と分母、
関係の分母を101で分解すると、最高次数12の既約因子33個が残る。
その外ではCramerの公式により経路基底と中心の五自由座標が特殊化できる。
各因子の通常な根でも両方の $i$ について全始点で経路・中心・四中心商を再計算した。
ゆえに全ての還元で中心が定階数となり、四切断が階数15の持上げに延びる。
前節と同じ自由性とKoszul性の判定が適用できる。
\end{proof}
